\documentclass[11pt]{article}

\usepackage[margin=1in]{geometry}
\usepackage{amsmath}
\usepackage{amssymb}
\usepackage{amsthm}
\usepackage{booktabs}
\usepackage{array}
\usepackage{enumitem}
\setlist{itemsep=2pt,parsep=2pt,topsep=3pt}
\usepackage[T1]{fontenc}
\usepackage[hidelinks]{hyperref}

\theoremstyle{plain}
\newtheorem{theorem}{Theorem}[section]
\newtheorem{lemma}[theorem]{Lemma}
\newtheorem{proposition}[theorem]{Proposition}
\newtheorem{corollary}[theorem]{Corollary}
\theoremstyle{definition}
\newtheorem{definition}[theorem]{Definition}
\theoremstyle{remark}
\newtheorem{remark}[theorem]{Remark}

\numberwithin{equation}{section}

\newcommand{\Sym}{\mathrm{S}}
\newcommand{\Alt}{\mathrm{A}}
\newcommand{\Gam}{\Gamma}
\DeclareMathOperator{\sgn}{sgn}
\DeclareMathOperator{\rank}{rank}
\newcommand{\normal}{\trianglelefteq}

\title{The parity subgroup of $\Sym_m \times \Sym_n$ is $2$-generated:\\
a proof of Fernandes' conjecture}

\author{Haoyu Chen\thanks{%
Disclosure of the production pipeline. The proof reproduced below was found by OpenAI GPT-5.6
running in the \texttt{codex} command-line interface, responding to a disguised problem
statement that named neither the conjecture nor its author. It was then verified line by line
by Claude (Fable 5), and submitted, in separate sessions and without the solver's reasoning
trace, to two adversarial reviewers: GPT-5.6 at ``Pro'' effort returned \textsc{valid} with no
issues, and Qwen3.8-Max returned \textsc{incomplete} with six presentation objections and no
critical error. Each of those six objections is answered in the present write-up; they are
listed in Section~\ref{sec:chain}. All small cases were checked by exhaustive enumeration
(Section~\ref{sec:comp}). Problem selection, task routing, review scheduling, and the drafting
of this paper were likewise automated; no human supplied a definition, a lemma, a proof idea,
or a correction. The argument is elementary and self-contained, and the reader is invited to
check it directly.}}

\date{16 August 2026}

\begin{document}

\maketitle

\begin{abstract}
For integers $m \ge n \ge 2$ let
$\Gam_{m,n} = \{(\sigma,\tau) \in \Sym_m \times \Sym_n : \sgn(\sigma) = \sgn(\tau)\}$
be the index-$2$ subgroup of pairs of permutations of equal parity. Fernandes
(arXiv:2605.12342, Conjecture~1) conjectured that $\Gam_{m,n}$ has rank $2$ for all
$m \ge n \ge 2$ outside the list $(m,n) \in \{(2,2),(3,3),(4,3),(4,4)\}$, and proved a special
case. We give an independent proof of the conjecture in full, by exhibiting a uniform pair of generators for every
admissible $(m,n)$. The proof is elementary: the generators are chosen so that both coordinate
projections are surjective, and Goursat's lemma then confines the generated subgroup to a fibre
product over a common quotient of $\Sym_m$ and $\Sym_n$. For unequal degrees the only such
quotient is $C_2$, which forces the fibre product to be $\Gam_{m,n}$ itself; for equal degrees
$r \ge 5$ the competing quotient $\Sym_r$ is killed by an order obstruction, the second
generator being chosen with components of orders $2$ and $r$ or $r-1$. The four exceptional
pairs are exactly those at which one of the two mechanisms fails, and we exhibit the failure
explicitly. All claims are accompanied by exhaustive machine verification for $m,n \le 7$,
including the closure computation for $|\Gam_{7,7}| = 12{,}700{,}800$. In addition, the main
theorem is now formally verified in Lean~4 against the \texttt{formal-conjectures} statement of
Conjecture~1, by a proof that is independent of the argument below rather than a transcription of
it (Section~\ref{sec:chain}). An earlier kernel-checked Lean~4 proof of the same statement was
submitted independently by Kenta Kitamura on 11~August~2026, predating this project; we do not
claim priority for the formal verification (see the Prior work paragraph in
Section~\ref{sec:chain}).
\end{abstract}

\section{Introduction}
\label{sec:intro}

\subsection{The conjecture}

Let $m,n \ge 2$ be integers, let $\Sym_m$ and $\Sym_n$ denote the symmetric groups on
$[m] = \{1,\dots,m\}$ and on a disjoint copy $[n'] = \{1',\dots,n'\}$, and set
\begin{equation}
\label{eq:gammadef}
\Gam_{m,n} \;=\; \bigl\{(\sigma,\tau) \in \Sym_m \times \Sym_n \;:\;
   \sgn(\sigma) = \sgn(\tau)\bigr\}.
\end{equation}
Equivalently, $\Gam_{m,n}$ is the kernel of the homomorphism
$\Sym_m \times \Sym_n \to \{\pm1\}$, $(\sigma,\tau) \mapsto \sgn(\sigma)\sgn(\tau)$; since
$m,n \ge 2$ this homomorphism is onto, so $\Gam_{m,n}$ has index $2$ and order $m!\,n!/2$. Under
the natural embedding of $\Sym_m \times \Sym_n$ into the symmetric group on
$[m] \cup [n']$, $\Gam_{m,n}$ is exactly the set of pairs whose joint action on
$[m] \cup [n']$ is an even permutation. Fernandes writes $\Gam_{m\oplus n}$ for this group; we
write $\Gam_{m,n}$, and we always assume $m \ge n$, which is no loss since
$\Gam_{m,n} \cong \Gam_{n,m}$.

The \emph{rank} of a finite group $G$, written $\rank(G)$, is the least cardinality of a
generating set. The group $\Gam_{m,n}$ arises in Fernandes' study of monoids of transformations
that are even on maximal proper subsets \cite{fernandes}, where its rank is computed for small
parameters and conjectured in general.

\begin{quote}
\textbf{Conjecture 1} (Fernandes, 2026 \cite{fernandes}). \emph{For all integers
$m \ge n \ge 2$ with $(m,n) \notin \{(2,2),(3,3),(4,3),(4,4)\}$, the group $\Gam_{m,n}$ has rank
$2$.}
\end{quote}

The four excluded pairs are genuine exceptions, not artefacts. Indeed
\[
\Gam_{2,2} = \bigl\{(\mathrm{id},\mathrm{id}),\; ((1\,2),(1'\,2'))\bigr\} \cong C_2
\]
has rank $1$, while $\Gam_{3,3}$, $\Gam_{4,3}$ and $\Gam_{4,4}$, of orders $18$, $72$ and $288$,
all have rank $3$. These values are recorded in \cite{fernandes} and were obtained there by
machine computation; we have reconfirmed all four by an independent exhaustive sweep over all
ordered pairs of elements (Section~\ref{sec:comp}).

Fernandes proves his conjecture in a special case, his Corollary~4, under a coprimality
hypothesis on the two degrees. The general case --- in particular, all the equal-degree cases
$m = n \ge 5$, where a coprimality hypothesis is unavailable --- was left open, and was
transcribed into Lean~4 as an open problem in the \texttt{formal-conjectures} corpus
\cite{formalconjectures} on 9 August 2026 (Issue \#4815, PR \#4836), together with four
\texttt{sorry}-ed variant statements recording the known ranks of the exceptional pairs.

Independently of the present project, and before it began, Kenta Kitamura (GitHub user
\texttt{KitaKen1}) opened \texttt{google-deepmind/formal-conjectures} pull request \#4868 on
11~August~2026, ``mark Fernandes conjecture 1 as solved,'' linking a kernel-checked Lean~4 proof
of the exact \texttt{formal-conjectures} statement of Conjecture~1; as of 18~August~2026 that pull
request remains open and unmerged. This predates the start of the work reported here
(16~August~2026). We learned of it only after completing our own proof and formalization, discuss
it further in the Prior work paragraph of Section~\ref{sec:chain}, and do not claim priority or
first resolution of the conjecture.

\subsection{Result}

\begin{theorem}
\label{thm:main}
Let $m \ge n \ge 2$ be integers with $(m,n) \notin \{(2,2),(3,3),(4,3),(4,4)\}$. Then
$\Gam_{m,n}$ is generated by two elements. Explicitly, with $a_r$, $b_r$, $d_r$ as defined in
\eqref{eq:abdef} and \eqref{eq:ddef} below:
\begin{enumerate}
\item if $m > n$, then $\Gam_{m,n} = \bigl\langle (a_m,a_n),\,(b_m,b_n) \bigr\rangle$;
\item if $m = n = r \ge 5$, then $\Gam_{r,r} = \bigl\langle (a_r,a_r),\,(b_r,d_r) \bigr\rangle$.
\end{enumerate}
Moreover $\rank(\Gam_{m,n}) = 2$ exactly.
\end{theorem}

Because $(4,3)$ is excluded, case (i) covers every admissible pair with $m > n$; case (ii)
covers every admissible pair with $m = n$, the pairs $(2,2)$, $(3,3)$ and $(4,4)$ being
excluded. So Theorem~\ref{thm:main} settles Conjecture~1. It subsumes Fernandes' Corollary~4:
the generators are uniform in $(m,n)$ and no arithmetic condition on the degrees is imposed.

\subsection{Method}

Write $\pi_1, \pi_2$ for the two coordinate projections of $\Sym_m \times \Sym_n$. A subgroup
$H$ is \emph{subdirect} if $\pi_1(H) = \Sym_m$ and $\pi_2(H) = \Sym_n$. The proof has three
moving parts.

\emph{(a) Surjectivity of both projections.} The elements $a_r$ (an even cycle) and $b_r$ (an odd
transposition) generate $\Sym_r$ for every $r \ge 2$, and the modified pair $a_r,d_r$ generates
as well because $b_r = d_r a_r^{-1}$. Hence both candidate subgroups are subdirect.

\emph{(b) Goursat's lemma.} A subdirect subgroup of $G \times K$ is a fibre product: there are
$N_G \normal G$, $N_K \normal K$ and an isomorphism $\phi : G/N_G \to K/N_K$ with
$H = \{(g,k) : \phi(gN_G) = kN_K\}$. The isomorphism type $Q$ of the common quotient $G/N_G
\cong K/N_K$ is thus an invariant of $H$, and $|H| = |G|\,|K|/|Q|$. If $Q$ is trivial then
$H$ is everything, which is impossible inside $\Gam_{m,n}$; and if $Q \cong C_2$ then the two
kernels are forced to be $\Alt_m$ and $\Alt_n$ and $H = \Gam_{m,n}$ exactly. So the entire
problem is to rule out every \emph{other} common quotient.

\emph{(c) Killing the remaining quotients.} The nontrivial quotients of $\Sym_r$ are $C_2$ and
$\Sym_r$ for $r \ne 4$, and $C_2$, $\Sym_3$, $\Sym_4$ for $r = 4$. For $m > n$ (excluding
$(4,3)$) an order count leaves $C_2$ as the only common quotient. For $m = n = r$ the quotient
$\Sym_r$ survives that count; it corresponds to $H$ being the graph of an automorphism $\phi$ of
$\Sym_r$. This is where the asymmetric second generator $(b_r,d_r)$ does its work: such a graph
would force $\phi(b_r) = d_r$, while $|b_r| = 2$ and $|d_r| \in \{r-1,r\}$. For $r \ge 5$ these
differ, and automorphisms preserve orders --- so nothing about $\mathrm{Aut}(\Sym_r)$, in
particular nothing about the outer automorphism of $\Sym_6$, needs to be known.

At $r = 3$ mechanism (c) fails because $|d_3| = 2 = |b_3|$; at $r = 4$, and at $(m,n) = (4,3)$,
it fails because of the extra quotient $\Sym_4/V_4 \cong \Sym_3$. Section~\ref{sec:exceptions}
makes both failures explicit, so that the shape of the exceptional list is explained rather than
merely respected.

The paper is organised as follows. Section~\ref{sec:gens} fixes notation and establishes the
properties of the generators. Sections~\ref{sec:goursat} and \ref{sec:normal} supply the two
group-theoretic inputs, with complete proofs. Section~\ref{sec:proof} proves
Theorem~\ref{thm:main}. Section~\ref{sec:exceptions} discusses the exceptional pairs,
Section~\ref{sec:comp} the computational verification, Section~\ref{sec:chain} the verification
chain, and Section~\ref{sec:limits} the limitations.

\section{Notation and the generating pairs}
\label{sec:gens}

Throughout, permutations act on the left and are composed \emph{right to left}: $(pq)(x) =
p(q(x))$. For $p \in \Sym_r$ and $i \ne j$ in $[r]$ we use repeatedly the conjugation rule
\begin{equation}
\label{eq:conj}
p\,(i\;j)\,p^{-1} \;=\; \bigl(p(i)\;p(j)\bigr),
\end{equation}
and, more generally, $p (i_1\,\cdots\,i_k) p^{-1} = (p(i_1)\,\cdots\,p(i_k))$. We write
$\Alt_r \normal \Sym_r$ for the alternating group, $\sgn : \Sym_r \to \{\pm1\}$ for the sign
homomorphism, $|g|$ for the order of a group element $g$, $V_4 \le \Sym_4$ for the Klein
four-group $\{\mathrm{id},(1\,2)(3\,4),(1\,3)(2\,4),(1\,4)(2\,3)\}$, and $C_2$ for the cyclic
group of order $2$. A $k$-cycle has sign $(-1)^{k-1}$; thus a cycle is an odd permutation
exactly when its length is even. The group $\Gam_{m,n}$ is as in \eqref{eq:gammadef}. Note for
later use that $\Alt_m \times \Alt_n \le \Gam_{m,n}$.

\begin{definition}
\label{def:ab}
For every integer $r \ge 2$ put
\begin{equation}
\label{eq:abdef}
a_r =
\begin{cases}
(1\;2\;\cdots\;r), & r \text{ odd},\\[2pt]
(1\;2\;\cdots\;r-1), & r \text{ even},
\end{cases}
\qquad
b_r =
\begin{cases}
(1\;2), & r \text{ odd},\\[2pt]
(r-1\;\;r), & r \text{ even},
\end{cases}
\end{equation}
with the convention that a displayed cycle of length $1$ denotes the identity; thus
$a_2 = \mathrm{id}$ and $b_2 = (1\,2)$. For $r \ge 3$ put
\begin{equation}
\label{eq:ddef}
d_r \;=\; b_r a_r .
\end{equation}
\end{definition}

\begin{lemma}
\label{lem:parity}
For every $r \ge 2$, the permutation $a_r$ is even and $b_r$ is odd.
\end{lemma}

\begin{proof}
If $r$ is odd, $a_r$ is an $r$-cycle of odd length, hence even, and $b_r$ is a transposition,
hence odd. If $r$ is even, $a_r$ is an $(r-1)$-cycle of odd length, hence even, and $b_r$ is
again a transposition. (For $r = 2$: $a_2 = \mathrm{id}$ is even and $b_2 = (1\,2)$ is odd.)
\end{proof}

\begin{lemma}
\label{lem:generate}
For every $r \ge 2$ we have $\langle a_r, b_r\rangle = \Sym_r$.
\end{lemma}

\begin{proof}
Suppose first that $r$ is odd, so $a := a_r = (1\,2\,\cdots\,r)$ and $b := b_r = (1\,2)$. By
\eqref{eq:conj}, for $0 \le k \le r-2$,
\[
a^k b a^{-k} \;=\; \bigl(a^k(1)\;\;a^k(2)\bigr) \;=\; (k+1\;\;k+2).
\]
So $\langle a,b\rangle$ contains all $r-1$ adjacent transpositions
$(1\,2), (2\,3), \dots, (r-1\;r)$. These generate $\Sym_r$: by induction on $j-i$, every
transposition $(i\;j)$ with $i < j$ lies in the group they generate, since
$(i\;\,j{+}1) = (j\;\,j{+}1)(i\;j)(j\;\,j{+}1)$ by \eqref{eq:conj}; and the transpositions
generate $\Sym_r$.

Now suppose $r$ is even, so $a := a_r = (1\,2\,\cdots\,r-1)$ fixes $r$, and
$b := b_r = (r-1\;\,r)$. By \eqref{eq:conj}, for $0 \le k \le r-2$,
\[
a^k b a^{-k} \;=\; \bigl(a^k(r-1)\;\;a^k(r)\bigr) \;=\; \bigl(a^k(r-1)\;\;r\bigr),
\]
and as $k$ runs through $0,1,\dots,r-2$ the point $a^k(r-1)$ runs through all of
$\{1,\dots,r-1\}$, because $a$ acts on that set as an $(r-1)$-cycle. Hence $\langle a,b\rangle$
contains every transposition $(i\;\,r)$ with $1 \le i \le r-1$; and then it contains every
transposition, since for distinct $i,j < r$ we have
$(i\;j) = (i\;\,r)(j\;\,r)(i\;\,r)$ by \eqref{eq:conj}. So $\langle a,b\rangle = \Sym_r$. This
argument includes $r = 2$, where $a_2 = \mathrm{id}$, the index $k$ takes only the value $0$,
and we obtain $(1\;2)$, which generates $\Sym_2$.
\end{proof}

\begin{lemma}
\label{lem:d}
Let $r \ge 3$ and $d_r = b_r a_r$. Then
\begin{equation}
\label{eq:dcycles}
d_r =
\begin{cases}
(2\;3\;\cdots\;r), & r \text{ odd},\\[2pt]
(1\;2\;\cdots\;r{-}2\;\;r\;\;r{-}1), & r \text{ even},
\end{cases}
\qquad\text{so}\qquad
|d_r| =
\begin{cases}
r-1, & r \text{ odd},\\
r, & r \text{ even}.
\end{cases}
\end{equation}
In both cases $d_r$ is an odd permutation, and $b_r = d_r a_r^{-1}$.
\end{lemma}

\begin{proof}
Recall $d_r(x) = b_r(a_r(x))$.

\emph{$r$ odd.} Here $a_r$ sends $x \mapsto x+1$ for $x < r$ and $r \mapsto 1$, while
$b_r = (1\,2)$. Then $d_r(1) = b_r(2) = 1$; for $2 \le x \le r-1$ we have
$d_r(x) = b_r(x+1) = x+1$, since $x+1 \ge 3$; and $d_r(r) = b_r(1) = 2$. Hence $d_r$ fixes $1$
and cycles $2 \mapsto 3 \mapsto \cdots \mapsto r \mapsto 2$, i.e.\ $d_r = (2\;3\;\cdots\;r)$, a
cycle of length $r-1$.

\emph{$r$ even} (so $r \ge 4$). Here $a_r$ sends $x \mapsto x+1$ for $x \le r-2$, sends
$r-1 \mapsto 1$ and fixes $r$; and $b_r = (r-1\;\,r)$. For $1 \le x \le r-3$ we get
$d_r(x) = b_r(x+1) = x+1$, since $2 \le x+1 \le r-2$; next $d_r(r-2) = b_r(r-1) = r$;
next $d_r(r-1) = b_r(1) = 1$, since $1 \notin \{r-1,r\}$ as $r \ge 4$; and finally
$d_r(r) = b_r(r) = r-1$. Hence $d_r$ is the single cycle
$1 \mapsto 2 \mapsto \cdots \mapsto r-2 \mapsto r \mapsto r-1 \mapsto 1$, that is
$d_r = (1\;2\;\cdots\;r{-}2\;\;r\;\;r{-}1)$, a cycle of length $r$.

The orders are the cycle lengths. For the parity: if $r$ is odd then $d_r$ is a cycle of even
length $r-1$, hence odd; if $r$ is even then $d_r$ is a cycle of even length $r$, hence odd. (Of
course this also follows from $\sgn(d_r) = \sgn(b_r)\sgn(a_r) = (-1)(+1)$ and
Lemma~\ref{lem:parity}; we have given the cycle computation because the orders are needed in any
case.) Finally $b_r = d_r a_r^{-1}$ is immediate from $d_r = b_r a_r$.
\end{proof}

\begin{lemma}
\label{lem:generate2}
For every $r \ge 3$ we have $\langle a_r, d_r \rangle = \Sym_r$.
\end{lemma}

\begin{proof}
By Lemma~\ref{lem:d}, $b_r = d_r a_r^{-1} \in \langle a_r,d_r\rangle$ and
$d_r = b_r a_r \in \langle a_r,b_r\rangle$, so $\langle a_r,d_r\rangle = \langle
a_r,b_r\rangle$, which is $\Sym_r$ by Lemma~\ref{lem:generate}.
\end{proof}

\section{Subdirect subgroups of a direct product}
\label{sec:goursat}

Let $G$ and $K$ be groups and let $\pi_G : G \times K \to G$ and $\pi_K : G\times K \to K$ be the
projections. A subgroup $H \le G \times K$ is \emph{subdirect} if $\pi_G(H) = G$ and
$\pi_K(H) = K$. The following is Goursat's lemma \cite{goursat} in the form we need. We give the
proof in full.

\begin{lemma}[Goursat]
\label{lem:goursat}
Let $H \le G \times K$ be subdirect. Put
\[
N_G = \{g \in G : (g,1) \in H\}, \qquad N_K = \{k \in K : (1,k) \in H\}.
\]
Then $N_G \normal G$ and $N_K \normal K$, and there is a (unique) isomorphism
$\phi : G/N_G \to K/N_K$ such that
\begin{equation}
\label{eq:fibre}
H \;=\; \bigl\{(g,k) \in G\times K \;:\; \phi(gN_G) = kN_K \bigr\}.
\end{equation}
In particular $|H| = |G|\,|K| / |Q|$, where $Q := G/N_G \cong K/N_K$.
\end{lemma}

\begin{proof}
\emph{$N_G$ and $N_K$ are normal.} That $N_G$ is a subgroup is clear from
$(g,1)(g',1)^{-1} = (g(g')^{-1},1)$. Let $g \in N_G$ and $x \in G$. Since $\pi_G(H) = G$, there
is $k \in K$ with $(x,k) \in H$; then
\[
(x,k)(g,1)(x,k)^{-1} = (xgx^{-1},\,1) \in H,
\]
so $xgx^{-1} \in N_G$. Hence $N_G \normal G$. Symmetrically $N_K \normal K$.

\emph{Construction of $\phi$.} Let $g \in G$. Since $\pi_G(H) = G$ there exists $k$ with
$(g,k) \in H$; we wish to set $\phi(gN_G) := kN_K$. We check this is unambiguous. Suppose
$(g,k) \in H$ and $(g',k') \in H$ with $gN_G = g'N_G$. Then
\[
(g,k)^{-1}(g',k') = (g^{-1}g',\, k^{-1}k') \in H,
\]
and $g^{-1}g' \in N_G$, so $(g^{-1}g',1) \in H$; multiplying,
\[
(g^{-1}g',1)^{-1}\,(g^{-1}g',\,k^{-1}k') = (1,\,k^{-1}k') \in H,
\]
whence $k^{-1}k' \in N_K$, i.e.\ $kN_K = k'N_K$. Taking $g = g'$ shows in particular that the
value does not depend on the choice of $k$; taking general $g,g'$ in the same coset shows it
depends only on $gN_G$. So $\phi : G/N_G \to K/N_K$ is a well-defined map.

\emph{$\phi$ is a homomorphism.} If $(g,k),(g',k') \in H$ then $(gg',kk') \in H$, so
$\phi(gg'N_G) = kk'N_K = \phi(gN_G)\phi(g'N_G)$.

\emph{$\phi$ is surjective.} Given $k \in K$, surjectivity of $\pi_K$ on $H$ gives $g$ with
$(g,k) \in H$, and then $\phi(gN_G) = kN_K$.

\emph{$\phi$ is injective.} Suppose $\phi(gN_G) = N_K$, witnessed by $(g,k) \in H$ with
$k \in N_K$. Then $(1,k) \in H$, hence $(g,1) = (g,k)(1,k)^{-1} \in H$, hence $g \in N_G$, i.e.\
$gN_G$ is trivial.

\emph{Proof of \eqref{eq:fibre}.} The inclusion $\subseteq$ is the definition of $\phi$.
Conversely, let $(g,k) \in G\times K$ satisfy $\phi(gN_G) = kN_K$. Choose $k_0$ with
$(g,k_0) \in H$; then $\phi(gN_G) = k_0N_K$, so $k_0N_K = kN_K$, so $k_0^{-1}k \in N_K$ and
$(1,k_0^{-1}k) \in H$. Therefore
\[
(g,k) = (g,k_0)\,(1,\,k_0^{-1}k) \in H .
\]

\emph{Uniqueness and order.} Any $\phi$ satisfying \eqref{eq:fibre} must send $gN_G$ to $kN_K$
whenever $(g,k) \in H$, which determines it. Finally, by \eqref{eq:fibre} the map
$H \to G$, $(g,k)\mapsto g$, is onto with all fibres of size $|N_K|$, so
$|H| = |G|\,|N_K| = |G|\,|K|/|Q|$.
\end{proof}

We refer to $Q = G/N_G \cong K/N_K$ as the \emph{common quotient} of the subdirect subgroup $H$.
Note the two degenerate cases: $Q = 1$ means $N_G = G$ and $N_K = K$, i.e.\ $H = G \times K$;
and $|Q| = |G| = |K|$ means $N_G = N_K = 1$, i.e.\ $H$ is the graph
$\{(g,\phi(g))\}$ of an isomorphism $\phi : G \to K$.

\section{Normal subgroups and quotients of symmetric groups}
\label{sec:normal}

\begin{lemma}
\label{lem:centralizer}
For $r \ge 4$ the centraliser of $\Alt_r$ in $\Sym_r$ is trivial. (For $r = 3$ it is
$\Alt_3 \cong C_3$, so the hypothesis $r \ge 4$ cannot be dropped.)
\end{lemma}

\begin{proof}
Let $r \ge 4$ and let $\sigma \in \Sym_r$ centralise $\Alt_r$, and suppose $\sigma \ne
\mathrm{id}$. Pick $i$ with $j := \sigma(i) \ne i$. Since $r \ge 4$ we may choose two further
points $k,l$ with $i,j,k,l$ pairwise distinct. The $3$-cycle $\theta = (i\;k\;l)$ lies in
$\Alt_r$, so by \eqref{eq:conj}
\[
\theta = \sigma\theta\sigma^{-1} = \bigl(\sigma(i)\;\;\sigma(k)\;\;\sigma(l)\bigr)
       = \bigl(j\;\;\sigma(k)\;\;\sigma(l)\bigr).
\]
Two $3$-cycles are equal only if they move the same three points, so $j \in \{i,k,l\}$,
contradicting the choice of $j,k,l$. Hence $\sigma = \mathrm{id}$. For $r = 3$, $\Alt_3$ is
abelian, so it centralises itself, while no transposition centralises $(1\,2\,3)$; the
centraliser is therefore $\Alt_3$.
\end{proof}

\begin{lemma}
\label{lem:s4v4}
$V_4 \normal \Sym_4$ and $\Sym_4/V_4 \cong \Sym_3$.
\end{lemma}

\begin{proof}
Let $\Pi$ be the set of the three partitions of $\{1,2,3,4\}$ into two unordered pairs:
\[
P_1 = \{\{1,2\},\{3,4\}\},\quad P_2 = \{\{1,3\},\{2,4\}\},\quad P_3 = \{\{1,4\},\{2,3\}\}.
\]
The natural action of $\Sym_4$ on subsets of $\{1,2,3,4\}$ permutes $\Pi$, giving a homomorphism
$\rho : \Sym_4 \to \mathrm{Sym}(\Pi) \cong \Sym_3$.

$\rho$ is surjective: $(1\,2)$ fixes $P_1$ and interchanges $P_2$ and $P_3$ (it sends $\{1,3\}$
to $\{2,3\}$ and $\{2,4\}$ to $\{1,4\}$, i.e.\ $P_2 \mapsto P_3$), while $(1\,3)$ fixes $P_2$ and
interchanges $P_1$ and $P_3$. So the image contains two distinct transpositions of $\Pi$, and
two distinct transpositions generate $\mathrm{Sym}(\Pi) \cong \Sym_3$.

$V_4 \le \ker\rho$: the element $(1\,2)(3\,4)$ fixes each of $\{1,2\}$ and $\{3,4\}$ setwise, so
fixes $P_1$; it sends $\{1,3\}$ to $\{2,4\}$ and $\{2,4\}$ to $\{1,3\}$, so fixes $P_2$; and it
sends $\{1,4\}$ to $\{2,3\}$ and $\{2,3\}$ to $\{1,4\}$, so fixes $P_3$. The same computation
applies to $(1\,3)(2\,4)$ and $(1\,4)(2\,3)$ after relabelling, and $\mathrm{id} \in \ker\rho$
trivially.

Since $\rho$ is surjective, $|\ker\rho| = |\Sym_4|/|\mathrm{Sym}(\Pi)| = 24/6 = 4 = |V_4|$, and
$V_4 \le \ker\rho$ forces $\ker\rho = V_4$. Being a kernel, $V_4 \normal \Sym_4$, and the first
isomorphism theorem gives $\Sym_4/V_4 \cong \mathrm{Sym}(\Pi) \cong \Sym_3$. (In particular the
quotient is nonabelian, so it is $\Sym_3$ and not $C_6$.)
\end{proof}

\begin{proposition}
\label{prop:normal}
The normal subgroups of $\Sym_r$ are:
\[
\begin{array}{c|l}
r & \text{normal subgroups of } \Sym_r\\ \hline
2 & 1,\ \Sym_2\\
3 & 1,\ \Alt_3,\ \Sym_3\\
4 & 1,\ V_4,\ \Alt_4,\ \Sym_4\\
r \ge 5 & 1,\ \Alt_r,\ \Sym_r
\end{array}
\]
Consequently the nontrivial quotients of $\Sym_r$, up to isomorphism, are
\begin{equation}
\label{eq:quotients}
\begin{array}{c|l}
r & \text{nontrivial quotients of } \Sym_r\\ \hline
2 & C_2\\
3 & \Sym_3,\ C_2\\
4 & \Sym_4,\ \Sym_3,\ C_2\\
r \ge 5 & \Sym_r,\ C_2
\end{array}
\end{equation}
and in every case the unique normal subgroup of index $2$ is $\Alt_r$.
\end{proposition}

\begin{proof}
$r = 2$: $\Sym_2 \cong C_2$ has only the two trivial subgroups, and $\Alt_2 = 1$.

$r = 3$: by Lagrange a proper nontrivial subgroup of $\Sym_3$ has order $2$ or $3$; the former
are the three subgroups generated by the transpositions, and the latter is $\Alt_3$. A subgroup
of order $2$ is not normal: the three transpositions are conjugate by \eqref{eq:conj}, so a
normal subgroup containing one contains all three and has order at least $4$, which does not
divide $6$. Thus the normal subgroups are $1, \Alt_3, \Sym_3$, and the nontrivial quotients are
$\Sym_3$ and $\Sym_3/\Alt_3 \cong C_2$.

$r = 4$: a normal subgroup is a union of conjugacy classes containing $\mathrm{id}$, and its
order divides $24$. The conjugacy classes of $\Sym_4$ are indexed by cycle type, with sizes
\[
1\ (\mathrm{id}),\quad 6\ \bigl(\text{type } 2\bigr),\quad 3\ \bigl(\text{type } 2{+}2\bigr),
\quad 8\ \bigl(\text{type }3\bigr),\quad 6\ \bigl(\text{type }4\bigr).
\]
Enumerating the sixteen sub-multisets of $\{6,3,8,6\}$ and adding $1$ for the identity class, the
attainable orders are
\[
1,\;4,\;7,\;9,\;10,\;12,\;13,\;15,\;16,\;18,\;21,\;24,
\]
and the only ones among them dividing $24$ are $1$, $4$, $12$, $24$.
These are realised only by $\{\mathrm{id}\}$; $\{\mathrm{id}\}\cup(\text{type }2{+}2) = V_4$;
$\{\mathrm{id}\}\cup(\text{type }2{+}2)\cup(\text{type }3) = \Alt_4$; and $\Sym_4$, all four of
which are indeed normal subgroups ($V_4$ by Lemma~\ref{lem:s4v4}, $\Alt_4$ as the kernel of
$\sgn$). The nontrivial quotients therefore have orders $24$, $6$, $2$, and are $\Sym_4$,
$\Sym_4/V_4 \cong \Sym_3$ by Lemma~\ref{lem:s4v4}, and $\Sym_4/\Alt_4 \cong C_2$.

$r \ge 5$: let $N \normal \Sym_r$. Then $N \cap \Alt_r \normal \Alt_r$, and $\Alt_r$ is simple
for $r \ge 5$, so $N \cap \Alt_r \in \{1, \Alt_r\}$. If $N \cap \Alt_r = \Alt_r$ then
$\Alt_r \le N \le \Sym_r$, and since $[\Sym_r : \Alt_r] = 2$ we get $N \in \{\Alt_r,\Sym_r\}$. If
$N \cap \Alt_r = 1$, then for $x \in N$ and $y \in \Alt_r$ the commutator
$[x,y] = x^{-1}(y^{-1}xy)$ lies in $N$ (as $N \normal \Sym_r$) and equals
$(x^{-1}y^{-1}x)y \in \Alt_r$ (as $\Alt_r \normal \Sym_r$), hence lies in $N \cap \Alt_r = 1$. So
$N$ centralises $\Alt_r$, and $N = 1$ by Lemma~\ref{lem:centralizer}. The quotients are then
$\Sym_r$ and $\Sym_r/\Alt_r \cong C_2$.

Finally, in each row of the table the only normal subgroup of index $2$ is $\Alt_r$: for $r = 2$
it is $1 = \Alt_2$; for $r = 3,4$ and $r \ge 5$ inspection of the lists (whose orders are
$1,3,6$; $1,4,12,24$; $1,r!/2,r!$) gives $\Alt_r$ in each case.
\end{proof}

The following consequence is the engine of the whole proof.

\begin{proposition}
\label{prop:key}
Let $m,n \ge 2$ and let $H \le \Gam_{m,n}$ be a subdirect subgroup of $\Sym_m \times \Sym_n$,
with common quotient $Q$ as in Lemma~\ref{lem:goursat}. Then:
\begin{enumerate}
\item $Q$ is nontrivial;
\item if $Q \cong C_2$ then $H = \Gam_{m,n}$.
\end{enumerate}
\end{proposition}

\begin{proof}
(i) If $Q$ were trivial then, as noted after Lemma~\ref{lem:goursat}, $H = \Sym_m \times \Sym_n$.
But $\Sym_m\times\Sym_n \not\le \Gam_{m,n}$: since $m \ge 2$ the transposition $(1\;2)$ lies in
$\Sym_m$, and the pair $\bigl((1\;2), \mathrm{id}\bigr)$ has
$\sgn((1\,2)) = -1 \ne +1 = \sgn(\mathrm{id})$, so it lies in $\Sym_m\times\Sym_n$ but not in
$\Gam_{m,n}$. This contradicts $H \le \Gam_{m,n}$.

(ii) Suppose $Q \cong C_2$. Then $N_{\Sym_m} \normal \Sym_m$ has index $2$, so
$N_{\Sym_m} = \Alt_m$ by Proposition~\ref{prop:normal}; likewise $N_{\Sym_n} = \Alt_n$. The map
$\sgn$ induces isomorphisms $\Sym_m/\Alt_m \xrightarrow{\ \sim\ } \{\pm1\}$ and
$\Sym_n/\Alt_n \xrightarrow{\ \sim\ } \{\pm1\}$, and $C_2$ has only the identity automorphism, so
under these identifications the Goursat isomorphism $\phi : \Sym_m/\Alt_m \to \Sym_n/\Alt_n$ is
the identity of $\{\pm 1\}$; that is, $\phi(\sigma \Alt_m) = \tau\Alt_n$ holds if and only if
$\sgn(\sigma) = \sgn(\tau)$. Now \eqref{eq:fibre} reads
\[
H = \bigl\{(\sigma,\tau) : \sgn(\sigma) = \sgn(\tau)\bigr\} = \Gam_{m,n}. \qedhere
\]
\end{proof}

\section{Proof of Theorem~\ref{thm:main}}
\label{sec:proof}

\subsection{Unequal degrees}

\begin{proposition}
\label{prop:unequal}
Let $m > n \ge 2$ with $(m,n) \ne (4,3)$. Then
$\Gam_{m,n} = \bigl\langle (a_m,a_n),\,(b_m,b_n) \bigr\rangle$.
\end{proposition}

\begin{proof}
Put $H = \langle (a_m,a_n), (b_m,b_n)\rangle$.

\emph{$H \le \Gam_{m,n}$.} By Lemma~\ref{lem:parity}, $a_m$ and $a_n$ are both even and $b_m$ and
$b_n$ are both odd, so both generators lie in $\Gam_{m,n}$, which is a group.

\emph{$H$ is subdirect.} The image of $H$ under the first projection contains $a_m$ and $b_m$,
hence is $\Sym_m$ by Lemma~\ref{lem:generate}; symmetrically the second projection is onto
$\Sym_n$.

\emph{The common quotient is $C_2$.} Let $Q$ be the common quotient of $H$ given by
Lemma~\ref{lem:goursat}; it is a quotient of $\Sym_m$ and also of $\Sym_n$, and it is nontrivial
by Proposition~\ref{prop:key}(i). We show $Q \cong C_2$, distinguishing three cases according to
$m$ (recall $m > n \ge 2$, so $m \ge 3$).

\emph{Case $m \ge 5$.} By \eqref{eq:quotients}, $Q \cong C_2$ or $Q \cong \Sym_m$. In the second
case $|Q| = m!$; but $Q$ is a quotient of $\Sym_n$, so $|Q|$ divides $n!$, whence $m! \le n!$,
contradicting $m > n$. So $Q \cong C_2$.

\emph{Case $m = 4$.} Then $n \in \{2,3\}$, and $n = 3$ is excluded by hypothesis, so $n = 2$. By
\eqref{eq:quotients} the only nontrivial quotient of $\Sym_2$ is $C_2$, so $Q \cong C_2$.

\emph{Case $m = 3$.} Then $n = 2$, and again $Q \cong C_2$.

By Proposition~\ref{prop:key}(ii), $H = \Gam_{m,n}$.
\end{proof}

\begin{remark}
The pair $(4,3)$ must be excluded here for a reason visible in the proof: $\Sym_4$ and $\Sym_3$
share the nontrivial quotient $\Sym_3$, by Lemma~\ref{lem:s4v4}. This is not merely a gap in the
argument --- Section~\ref{sec:exceptions} shows that the displayed pair really does fail to
generate $\Gam_{4,3}$, and that it fails precisely by landing in the fibre product over $\Sym_3$.
\end{remark}

\subsection{Equal degrees}

\begin{proposition}
\label{prop:equal}
Let $r \ge 5$. Then $\Gam_{r,r} = \bigl\langle (a_r,a_r),\,(b_r,d_r)\bigr\rangle$.
\end{proposition}

\begin{proof}
Put $H = \langle (a_r,a_r), (b_r,d_r)\rangle$.

\emph{$H \le \Gam_{r,r}$.} By Lemma~\ref{lem:parity} $a_r$ is even, so $(a_r,a_r) \in
\Gam_{r,r}$; by Lemma~\ref{lem:parity} $b_r$ is odd, and by Lemma~\ref{lem:d} $d_r$ is odd as
well, so $(b_r,d_r) \in \Gam_{r,r}$.

\emph{$H$ is subdirect.} The first projection of $H$ contains $a_r$ and $b_r$, so is $\Sym_r$ by
Lemma~\ref{lem:generate}. The second projection contains $a_r$ and $d_r$, so is $\Sym_r$ by
Lemma~\ref{lem:generate2}.

\emph{The common quotient is $C_2$.} Let $Q$ be the common quotient. It is nontrivial by
Proposition~\ref{prop:key}(i), and since $r \ge 5$, \eqref{eq:quotients} leaves
$Q \cong C_2$ or $Q \cong \Sym_r$. Suppose $Q \cong \Sym_r$. Then
$|Q| = r! = |\Sym_r|$ forces both Goursat kernels to be trivial, so by the remark following
Lemma~\ref{lem:goursat} the subgroup $H$ is the graph of an automorphism $\phi$ of $\Sym_r$:
\[
H = \{(\sigma, \phi(\sigma)) : \sigma \in \Sym_r\}.
\]
Both generators of $H$ have this form, so
\[
\phi(a_r) = a_r, \qquad \phi(b_r) = d_r .
\]
But an automorphism preserves the order of every element, while $|b_r| = 2$ and, by
Lemma~\ref{lem:d},
\[
|d_r| = \begin{cases} r-1 \ge 4, & r \text{ odd},\\ r \ge 6, & r \text{ even},\end{cases}
\]
using $r \ge 5$ in both branches. Hence $|b_r| \ne |d_r|$, a contradiction. Therefore
$Q \cong C_2$, and Proposition~\ref{prop:key}(ii) gives $H = \Gam_{r,r}$.
\end{proof}

\begin{remark}
\label{rem:s6}
The obstruction just used is an order obstruction, and orders are preserved by \emph{every}
automorphism, inner or outer. In particular the argument requires no knowledge of
$\mathrm{Aut}(\Sym_r)$; the well-known exceptional outer automorphism of $\Sym_6$ is immaterial,
and the case $r = 6$ needs no separate treatment.
\end{remark}

\subsection{Conclusion, and the exact value of the rank}

\begin{proof}[Proof of Theorem~\ref{thm:main}]
Let $m \ge n \ge 2$ with $(m,n) \notin \{(2,2),(3,3),(4,3),(4,4)\}$. If $m > n$ then
$(m,n) \ne (4,3)$ and Proposition~\ref{prop:unequal} applies. If $m = n = r$ then
$r \notin\{2,3,4\}$, i.e.\ $r \ge 5$, and Proposition~\ref{prop:equal} applies. In both cases
$\Gam_{m,n}$ is generated by the two displayed elements, so $\rank(\Gam_{m,n}) \le 2$. The
reverse inequality is Corollary~\ref{cor:rank} below.
\end{proof}

\begin{corollary}
\label{cor:rank}
For every $(m,n)$ as in Theorem~\ref{thm:main}, the group $\Gam_{m,n}$ is nonabelian; hence it is
not cyclic, $\rank(\Gam_{m,n}) \ge 2$, and therefore $\rank(\Gam_{m,n}) = 2$.
\end{corollary}

\begin{proof}
A group of rank at most $1$ is cyclic, hence abelian, so it suffices to prove that $\Gam_{m,n}$
is nonabelian. Note $m \ge 3$: the only admissible pair with $m = 2$ would be $(2,2)$, which is
excluded.

If $m \ge 4$, then $\Alt_m \times \{\mathrm{id}\} \le \Alt_m\times\Alt_n \le \Gam_{m,n}$ and
$\Alt_m$ is nonabelian, since $(1\,2\,3)(1\,2\,4) = (1\,3)(2\,4)$ while
$(1\,2\,4)(1\,2\,3) = (1\,4)(2\,3)$.

If $m = 3$, then $n = 2$, because $(3,3)$ is excluded and $n \le m$. The first projection
$\Gam_{3,2} \to \Sym_3$ is surjective (given $\sigma \in \Sym_3$, pair it with the element of
$\Sym_2$ of the same sign) and injective (its kernel is $\{(\mathrm{id},\tau) : \sgn(\tau) = +1\}
= \{(\mathrm{id},\mathrm{id})\}$), so $\Gam_{3,2} \cong \Sym_3$, which is nonabelian.
\end{proof}

Together with the known values $\rank(\Gam_{2,2}) = 1$ and
$\rank(\Gam_{3,3}) = \rank(\Gam_{4,3}) = \rank(\Gam_{4,4}) = 3$, Theorem~\ref{thm:main}
determines $\rank(\Gam_{m,n})$ for all $m \ge n \ge 2$.

\section{The exceptional pairs}
\label{sec:exceptions}

The exceptional list in Conjecture~1 is exactly the list of pairs at which one of the two
mechanisms of Section~\ref{sec:proof} breaks down, and the breakdown is visible in the
constructions themselves; this explains, rather than merely accommodates, the shape of the
exceptional set. Orders quoted here were computed by exhaustive closure
(Section~\ref{sec:comp}).

\emph{$(m,n) = (4,3)$.} The unequal-degree pair of Proposition~\ref{prop:unequal} reads
$\bigl((1\,2\,3),(1\,2\,3)\bigr)$, $\bigl((3\,4),(1\,2)\bigr)$. It generates a subgroup of order
$24$ inside $\Gam_{4,3}$, which has order $72$; the index is $3$. This is exactly the fibre
product predicted by Lemma~\ref{lem:goursat} over the common quotient $\Sym_3$, which is a
quotient of $\Sym_4$ (namely $\Sym_4/V_4$, by Lemma~\ref{lem:s4v4}) as well as of $\Sym_3$:
indeed $|\Sym_4|\,|\Sym_3|/|\Sym_3| = 24$.

\emph{$(m,n) = (4,4)$.} The equal-degree pair reads $\bigl((1\,2\,3),(1\,2\,3)\bigr)$,
$\bigl((3\,4),(1\,2\,4\,3)\bigr)$; here the order obstruction of Proposition~\ref{prop:equal} is
available ($|b_4| = 2 \ne 4 = |d_4|$) and does rule out the quotient $\Sym_4$, but the extra
quotient $\Sym_4/V_4\cong\Sym_3$ is not ruled out --- and it is precisely what occurs: the
generated subgroup has order $96 = 24\cdot 24/6$, of index $3$ in $|\Gam_{4,4}| = 288$.

\emph{$(m,n) = (3,3)$.} Here $|d_3| = |(2\,3)| = 2 = |b_3|$, so the order obstruction is vacuous,
and the graph case really occurs. Conjugation by $a_3 = (1\,2\,3)$ is an automorphism $\phi$ of
$\Sym_3$ with $\phi(a_3) = a_3$ and, by \eqref{eq:conj},
$\phi(b_3) = \phi((1\,2)) = (a_3(1)\;a_3(2)) = (2\,3) = d_3$. Hence the subgroup generated by
$(a_3,a_3)$ and $(b_3,d_3)$ is contained in --- and, both projections being onto, equal to --- the
graph $\{(\sigma, a_3\sigma a_3^{-1})\}$, of order $6$ and index $3$ in $|\Gam_{3,3}| = 18$.

\emph{$(m,n) = (2,2)$.} $\Gam_{2,2} \cong C_2$ is cyclic of rank $1$, so it is excluded from a
conjecture asserting rank exactly $2$, though it is of course $2$-generated in the weak sense.

In all three nontrivial cases the recipe fails by exactly the index $3$ predicted by the
corresponding Goursat quotient of order $6$. That no \emph{other} choice of two elements
succeeds either --- the assertion $\rank = 3$ recorded in \cite{fernandes} --- we have
reconfirmed independently in Section~\ref{sec:comp}.

\section{Computational verification}
\label{sec:comp}

The proof of Theorem~\ref{thm:main} is complete and uses no computation. Nevertheless, since the
theorem asserts something checkable for each small $(m,n)$, we verified the explicit generating
pairs by exhaustive closure. The computation is a breadth-first search in the Cayley graph:
starting from the identity of $\Sym_m\times\Sym_n$ and repeatedly left-multiplying by the two
generators, one enumerates the generated subgroup and compares its cardinality with
$|\Gam_{m,n}| = m!\,n!/2$. Permutations were indexed in lexicographic order and the two
generators were pre-tabulated as index maps, so that each edge of the search costs two array
lookups. All arithmetic is exact.

\begin{table}[tb]
\centering
\small
\begin{tabular}{llrrl}
\toprule
$(m,n)$ & recipe & $|\langle\,\cdot\,\rangle|$ & $|\Gam_{m,n}| = m!n!/2$ & verdict\\
\midrule
$(3,2)$ & unequal & $6$ & $6$ & equal\\
$(4,2)$ & unequal & $24$ & $24$ & equal\\
$(5,2)$ & unequal & $120$ & $120$ & equal\\
$(5,3)$ & unequal & $360$ & $360$ & equal\\
$(5,4)$ & unequal & $1{,}440$ & $1{,}440$ & equal\\
$(6,2)$ & unequal & $720$ & $720$ & equal\\
$(6,3)$ & unequal & $2{,}160$ & $2{,}160$ & equal\\
$(6,4)$ & unequal & $8{,}640$ & $8{,}640$ & equal\\
$(6,5)$ & unequal & $43{,}200$ & $43{,}200$ & equal\\
$(7,3)$ & unequal & $15{,}120$ & $15{,}120$ & equal\\
$(7,6)$ & unequal & $1{,}814{,}400$ & $1{,}814{,}400$ & equal\\
$(5,5)$ & equal & $7{,}200$ & $7{,}200$ & equal\\
$(6,6)$ & equal & $259{,}200$ & $259{,}200$ & equal\\
$(7,7)$ & equal & $12{,}700{,}800$ & $12{,}700{,}800$ & equal\\
\midrule
$(4,3)$ & unequal & $24$ & $72$ & proper, index $3$\\
$(3,3)$ & equal & $6$ & $18$ & proper, index $3$\\
$(4,4)$ & equal & $96$ & $288$ & proper, index $3$\\
\bottomrule
\end{tabular}
\caption{Exhaustive closure of the generating pairs of Theorem~\ref{thm:main}. The first
fourteen rows are admissible pairs and the closure is all of $\Gam_{m,n}$. The last three rows
are the nontrivial exceptional pairs, where the recipe demonstrably fails, by the index
predicted in Section~\ref{sec:exceptions}.}
\label{tab:closure}
\end{table}

Table~\ref{tab:closure} reports the outcome for fourteen admissible pairs, all those with
$m \le 7$ that we selected, together with the three nontrivial exceptional pairs. In every
admissible case the closure has exactly $m!\,n!/2$ elements, the largest being
$|\Gam_{7,7}| = 12{,}700{,}800$; in every exceptional case it is a proper subgroup of index $3$.
Four further families of checks were run.

\begin{enumerate}
\item \emph{Generators.} For $2 \le r \le 12$: $\sgn(a_r) = +1$, $\sgn(b_r) = \sgn(d_r) = -1$,
$|b_r| = 2$, $|d_r| = r-1$ for odd $r$ and $r$ for even $r$, the closed form \eqref{eq:dcycles}
for $d_r$, and the identity $b_r = d_ra_r^{-1}$: all confirmed. For $2 \le r \le 8$,
$|\langle a_r,b_r\rangle| = |\langle a_r,d_r\rangle| = r!$: confirmed
(Lemmas~\ref{lem:generate}, \ref{lem:generate2}).

\item \emph{Ranks of the exceptional groups.} By exhaustive sweep over all ordered pairs of
elements: $\Gam_{2,2}$ (order $2$) is cyclic; and for $\Gam_{3,3}$ (order $18$), $\Gam_{4,3}$
(order $72$) and $\Gam_{4,4}$ (order $288$) \emph{no} pair of elements generates. This
independently reconfirms the values $\rank = 3$ quoted from \cite{fernandes} and used in
Section~\ref{sec:exceptions}, and it is what makes the exceptional list genuinely exceptional
rather than an artefact of our particular recipe.

\item \emph{Structural lemmas.} The conjugacy-class sizes of $\Sym_4$ were recomputed
($1,6,3,8,6$ for types $\mathrm{id}$, $2$, $2{+}2$, $3$, $4$), confirming the enumeration in
Proposition~\ref{prop:normal}; the homomorphism $\rho:\Sym_4\to\mathrm{Sym}(\Pi)$ of
Lemma~\ref{lem:s4v4} was constructed explicitly and found to have image of order $6$, nonabelian,
and kernel exactly $V_4$; and the centraliser of $\Alt_r$ in $\Sym_r$ was computed for
$3 \le r \le 6$, giving order $3$ for $r = 3$ and order $1$ for $r = 4,5,6$, confirming both
Lemma~\ref{lem:centralizer} and the necessity of its hypothesis $r \ge 4$.

\item \emph{The isomorphism $\Gam_{3,2}\cong\Sym_3$} used in Corollary~\ref{cor:rank} was
confirmed by checking that the first projection is injective on $\Gam_{3,2}$, which has order
$6$.
\end{enumerate}

None of these computations is load-bearing for Theorem~\ref{thm:main}; all of them agree with
it.

\section{The verification chain}
\label{sec:chain}

We describe how this result was produced and checked, because the answer bears on how much
weight a reader should give it before checking it personally.

\emph{Solver/verifier separation.} The proof was found by OpenAI GPT-5.6 running in the
\texttt{codex} command-line interface. The problem was presented in disguised form: the task
file stated the group $\Gam_{m,n}$, the exceptional list and the known ranks, but named neither
Fernandes nor the arXiv identifier, and the solver was instructed not to search the internet and
not to read any other file in the workspace. It was required to present only what it could
rigorously prove, and was explicitly permitted to return a partial answer; the solution it
returned covered every admissible pair and was correct.

\emph{Line-by-line verification.} The argument was then checked step by step by Claude
(Fable 5), including independent recomputation of the cycle forms, parities and orders of
$a_r,b_r,d_r$, of the conjugacy-class enumeration for $\Sym_4$, and of the closures of
Section~\ref{sec:comp}.

\emph{Adversarial review.} The argument was submitted, in two separate sessions with no access
to the solver's reasoning trace, to two reviewers, under a protocol that requires each objection
to be classified as a \emph{critical error} (breaking the logical chain) or a \emph{justification
gap} (a true step insufficiently argued), and requires a final verdict. One reviewer was
Qwen3.8-Max, from a different vendor than the solver; the other was GPT-5.6 at ``Pro'' effort,
which is the same vendor as the solver in a different harness and at a different effort setting,
and its independence should therefore be discounted accordingly.

GPT-5.6 returned \textsc{valid}, with no issues, having independently confirmed the cycle form of
$b_ra_r$ under right-to-left composition and the resulting orders $r-1$ and $r$, the generation
statement including the edge case $r = 2$, the normal-quotient classification, the exclusion of
quotients other than $C_2$ in the unequal case, the surjectivity of the second projection via
$b_r = d_ra_r^{-1}$, and the order argument against the graph case including $r = 6$.

Qwen3.8-Max returned \textsc{incomplete}: no critical errors, but six justification gaps, all of
them about what was asserted rather than proved. We list each objection with the place in this
write-up that answers it; we agree with the reviewer that all six were gaps in exposition and
none was a gap in the mathematics.

\begin{enumerate}
\item \emph{$\Gam_{m,n}$ was used before being defined.} Answered by \eqref{eq:gammadef}.

\item \emph{The well-definedness, bijectivity and converse inclusion in the Goursat lemma were
asserted, not proved.} Answered by the complete proof of Lemma~\ref{lem:goursat}.

\item \emph{$\Sym_4/V_4 \cong \Sym_3$ was asserted, whereas the class-size argument bounds only
the order of the quotient, leaving $C_6$ formally open.} Answered by Lemma~\ref{lem:s4v4}, which
exhibits the action of $\Sym_4$ on the three pair-partitions of $\{1,2,3,4\}$ and identifies the
kernel, and thereby also shows the quotient nonabelian.

\item \emph{The triviality of the centraliser of $\Alt_r$ in $\Sym_r$ was used without proof.}
Answered by Lemma~\ref{lem:centralizer}, which also records that its hypothesis $r \ge 4$ cannot
be dropped; it is applied only for $r \ge 5$, the cases $r \le 4$ of
Proposition~\ref{prop:normal} being handled by direct enumeration.

\item \emph{The parity of $d_r$ was not derived.} Answered by Lemma~\ref{lem:d}, which computes
the cycle form of $d_r$ and reads the parity off the (even) cycle length in both branches.

\item \emph{The step ``a trivial common quotient would make $H = \Sym_m\times\Sym_n$, which is
not contained in $\Gam_{m,n}$'' did not say why the containment fails.} Answered by the explicit
witness $((1\,2),\mathrm{id})$ in the proof of Proposition~\ref{prop:key}(i).
\end{enumerate}

\emph{Computational verification.} As reported in Section~\ref{sec:comp}.

\emph{Formal verification.} The main theorem is now formally verified. The Lean~4 file
\texttt{Fernandes.lean} (507 lines) compiles with zero errors and zero warnings, contains no
\texttt{sorry} and no \texttt{native\_decide}, and introduces no new axioms: \texttt{\#print
axioms} on each of its five public theorems reports exactly \texttt{[propext, Classical.choice,
Quot.sound]}, the three axioms in routine use throughout mathlib. The statement of the main
theorem, \texttt{conjecture\_1} --- including the auxiliary definitions \texttt{signDiffHom} and
\texttt{gammaSubgroup} --- is byte-identical to the statement of \texttt{conjecture\_1} in
\texttt{FormalConjectures/Arxiv/2605.12342/Conjecture1.lean} \cite{formalconjectures}; only the
proof body differs. We also formalised the variant
\texttt{conjecture\_1.variants.rank\_2\_2} ($\Gam_{2,2}$ is cyclic), discharging a second of the
five \texttt{sorry} placeholders in that file. The remaining three variants, recording that
$\Gam_{3,3}$, $\Gam_{4,3}$ and $\Gam_{4,4}$ have rank $3$, are not formalised; their verification
remains the computational argument of Section~\ref{sec:comp}.

\emph{Formalisation notes.} The formal proof is not a line-by-line translation of
Sections~\ref{sec:gens}--\ref{sec:proof}, but an independently verified route to the same
statement, which is arguably additional evidence for the theorem rather than a mechanical
transcription of the write-up above. Four simplifications are worth recording. First, mathlib's
\texttt{Equiv.Perm.alternatingGroup\_le\_of\_normal} (for $r \ge 5$) replaces
Lemma~\ref{lem:centralizer} and Proposition~\ref{prop:normal} entirely. Second, the formalisation
uses a single generating pair uniformly in $r$ --- $c = \mathrm{finRotate}\,r$ and
$t = \mathrm{swap}\,0\,1$ --- whose roles under the sign character swap automatically with the
parity of $r$, so the even/odd case split of Section~\ref{sec:gens} never appears. Third, the
Goursat step is carried out with kernels alone (\texttt{goursatFst}, \texttt{goursatSnd}); the
quotient isomorphism $\phi$ of Lemma~\ref{lem:goursat} is never constructed. Fourth, the order
obstruction of Proposition~\ref{prop:equal} is weakened, and suffices in this weaker form, to
$d^2 \ne 1$.

\emph{Prior work.} Kenta Kitamura (GitHub user \texttt{KitaKen1}) has an earlier, independent
kernel-checked Lean~4 proof of the exact \texttt{formal-conjectures} statement of Conjecture~1,
linked from \texttt{google-deepmind/formal-conjectures} pull request \#4868 (``mark Fernandes
conjecture 1 as solved''), opened 11~August~2026 and, as of this writing (18~August~2026), still
open and unmerged. That date precedes the start of the present project (16~August~2026), so
Kitamura's formalization has priority; our Lean development described above and the
elementary argument of Sections~\ref{sec:gens}--\ref{sec:proof} were both found independently,
by a different, unrelated pipeline, without knowledge of his proof. We record our proof anyway
because we believe the human-readable argument here --- via Goursat's lemma and the classification
of normal subgroups of $\Sym_r$ --- may retain expository value alongside the formal one, even
though it establishes nothing that was not already machine-checked.

\section{Limitations}
\label{sec:limits}

\begin{enumerate}
\item \emph{The proof is elementary.} It combines two textbook ingredients --- Goursat's lemma
and the classification of normal subgroups of $\Sym_r$ --- with one small design choice, namely
replacing the symmetric second generator $(b_r,b_r)$ by $(b_r,d_r)$ so that the two components
have different orders. There is no new technique here. The value of the result is that it
settles a conjecture stated in 2026 and open at the time of writing, not that the method is
novel.

\item \emph{The equal-degree construction is not canonical.} Many other second generators would
work; ours is chosen to make the order obstruction as short as possible. We have not
investigated which pairs of elements generate $\Gam_{r,r}$ in general, nor the proportion of
generating pairs (a natural question in the tradition of P.~Hall's Eulerian functions
\cite{hall}).

\item \emph{Literature work was automated, and we did not consult the source article.} Our
statement of Conjecture~1, of the known ranks of the four exceptional pairs, and of the scope of
Fernandes' Corollary~4, is taken from the Lean transcription in \texttt{formal-conjectures}
\cite{formalconjectures} and from an automated literature pass, not from the arXiv PDF; the same
applies to \cite{fv}, cited only as background. Theorem~\ref{thm:main} and its proof are
self-contained and independent of all of that, but attributions --- in particular our
characterisation of what Corollary~4 of \cite{fernandes} covers --- should be read with the
caveat. An initial automated search for a prior proof of Conjecture~1 found none; a later, more
targeted search did locate one --- K.~Kitamura's independent Lean~4 formalization, predating this
project (Section~\ref{sec:chain}) --- so we no longer treat the earlier negative search as
informative, and any remaining priority claim on our part should be discounted accordingly.

\item \emph{Formalisation is complete for the main theorem and one variant.} The Lean~4
formalisation (Section~\ref{sec:chain}) proves Theorem~\ref{thm:main} against the
\texttt{formal-conjectures} statement of Conjecture~1, together with the variant recording
$\rank(\Gam_{2,2}) = 1$. The three remaining variants, for the exceptional ranks
$\rank(\Gam_{3,3}) = \rank(\Gam_{4,3}) = \rank(\Gam_{4,4}) = 3$, are not formalised; their
verification remains the exhaustive computation of Section~\ref{sec:comp}, which is a complete
proof but not a machine-checked one in the Lean sense.

\item \emph{The exceptional cases are quoted, not proved here.} We reconfirmed
$\rank(\Gam_{3,3}) = \rank(\Gam_{4,3}) = \rank(\Gam_{4,4}) = 3$ by exhaustive computation, which
is a complete proof for those three finite groups, but a computational one; we give no
structural explanation of why no pair generates. Conjecture~1 as stated presupposes those
values, though Theorem~\ref{thm:main} does not depend on them: the assertion
$\rank(\Gam_{m,n}) = 2$ for admissible $(m,n)$ is proved in Corollary~\ref{cor:rank} without
reference to them.
\end{enumerate}

\subsection*{Acknowledgement of process}

No human contributed a mathematical step to this paper. The pipeline --- problem selection,
disguised task construction, solving, verification, adversarial review, computation, and
drafting --- was run autonomously. Responsibility for the correctness of the argument
nevertheless rests where it always does: with the reader who checks it.

\begin{thebibliography}{9}

\bibitem{fernandes}
V.~H. Fernandes,
\emph{Groups of permutations that are even on maximal proper subsets, and related monoids},
arXiv:2605.12342 (2026).
Conjecture~1 is the statement proved here; Corollary~4 of that paper establishes a special case.

\bibitem{fv}
V.~H. Fernandes and A.~Vernitski,
\emph{Internat. J. Algebra Comput.} (2026).
Cited as background on the ranks of the monoids in whose study $\Gam_{m,n}$ arises.
Bibliographic details are as supplied by our automated literature pass; we did not consult the
article (see Section~\ref{sec:limits}).

\bibitem{goursat}
\'E. Goursat,
\emph{Sur les substitutions orthogonales et les divisions r\'eguli\`eres de l'espace},
Ann. Sci. \'Ecole Norm. Sup. (3) \textbf{6} (1889), 9--102.
The source of Lemma~\ref{lem:goursat}.

\bibitem{hall}
P.~Hall,
\emph{The Eulerian functions of a group},
Quart. J. Math. Oxford Ser. \textbf{7} (1936), 134--151.

\bibitem{formalconjectures}
Google DeepMind and contributors,
\emph{formal-conjectures}: a corpus of open conjectures stated in Lean~4,
\url{https://github.com/google-deepmind/formal-conjectures}.
The statement proved here was added as
\texttt{FormalConjectures/Arxiv/2605.12342/Conjecture1.lean} on 9~August~2026
(Issue \#4815, PR \#4836), with five \texttt{sorry} placeholders: the conjecture itself and
four variants recording the ranks of the exceptional pairs.

\bibitem{mathlib}
The mathlib Community,
\emph{The Lean mathematical library},
Proc. 9th ACM SIGPLAN Int. Conf. on Certified Programs and Proofs (CPP 2020), 367--381.

\end{thebibliography}

\end{document}
